% =====================================================================
% Manuscript/macros.tex
% Canonical TeX preamble macros for the unified manuscript.
% Produced by Agent 17 (Notation and Normalization Unification, 2026-05-13).
% This file is to be \input by Manuscript/main.tex (Agent 19).
%
% The macros are taken verbatim from the canonical Notation Register
% §15 (Manuscript/00_notation_register.md). No draft introduces its
% own \newcommand definitions; every macro used below appears in some
% draft and is consistent across drafts.
%
% Package requirements (loaded in main.tex):
%   amsmath, amssymb, mathtools (for \DeclareMathOperator and
%   common math symbols).
% =====================================================================

% --- Ambient space and basic symbols ---
\newcommand{\R}{\mathbb{R}}
\newcommand{\Om}{\Omega}
\newcommand{\Otilde}{\widetilde{\Omega}}
\newcommand{\eps}{\varepsilon}

% --- Asymmetry, balls, deficits ---
\newcommand{\Fr}{\mathcal{A}}                  % canonical Fraenkel asymmetry symbol
\newcommand{\Bstar}{B^{*}}                     % equal-volume ball at the origin
\newcommand{\deltaT}{\delta_{T}}               % Saint-Venant deficit
\newcommand{\deltaFK}{\delta_{\mathrm{FK}}}    % Faber-Krahn deficit
\newcommand{\BDValpha}{\alpha_{\mathrm{BDV}}}  % BDV smooth asymmetry

% --- Foliation radii and centre fields ---
\newcommand{\rstar}{\rho_{*}}
\newcommand{\rhad}{{\widehat{\rho}}}
\newcommand{\rhodelta}{\rho_{\delta}}
\newcommand{\zbar}{\bar z_{\rho}}              % bulk centroid of E_rho

% --- Level-set foliation and velocities ---
\newcommand{\Erho}{E_{\rho}}
\newcommand{\Frho}{F_{\rho}}
\newcommand{\Vrho}{V_{\rho}}
\newcommand{\nurho}{\nu_{\rho}}

% --- Metric quotient space ---
\newcommand{\X}{\mathcal{X}}
\newcommand{\dX}{d_{\mathcal{X}}}

% --- Pi block (auxiliary, Agent 14) ---
\newcommand{\PiE}{\Pi(E_{\rho},\zbar)}
\newcommand{\OmPi}{\Omega_{\Pi}}

% --- Operators ---
\DeclareMathOperator{\dist}{dist}
\DeclareMathOperator{\Div}{div}
\DeclareMathOperator{\Lip}{Lip}
\DeclareMathOperator{\diam}{diam}
\providecommand{\hH}{\mathcal{H}}
\providecommand{\NS}{\mathrm{NS}}
\providecommand{\fint}{\Xint-}
\def\Xint#1{\mathchoice{\XXint\displaystyle\textstyle{#1}}%
   {\XXint\textstyle\scriptstyle{#1}}%
   {\XXint\scriptstyle\scriptscriptstyle{#1}}%
   {\XXint\scriptscriptstyle\scriptscriptstyle{#1}}%
   \!\int}
\def\XXint#1#2#3{{\setbox0=\hbox{$#1{#2#3}{\int}$}
     \vcenter{\hbox{$#2#3$}}\kern-.5\wd0}}
